\section{The FunQA Dataset}
\begin{figure}[!t]
    \centering
    % \begin{subfigure}[b]{0.3\textwidth}
    %     \centering
    %     \includegraphics[height=3cm]{figures/fig_word_cloud_h.png}
    %     \caption{Word cloud of HumorQA}
    %     \label{fig:wc_h}
    % \end{subfigure}
    % \hfill

    % \begin{subfigure}[b]{0.3\textwidth}
    %     \centering
    %     \includegraphics[height=3cm]{figures/fig_word_cloud_c.png}
    %     \caption{Word cloud of CreativeQA}
    %     \label{fig:wc_c}
    % \end{subfigure}
    % \hfill

    % \begin{subfigure}[b]{0.3\textwidth}
    %     \centering
    %     \includegraphics[height=3cm]{figures/fig_word_cloud_m.png}
    %     \caption{Word cloud of MagicQA}
    %     \label{fig:wc_m}
    % \end{subfigure}

    \begin{subfigure}[b]{0.45\textwidth}
        \centering
        \includegraphics[width=1\textwidth]
        {figures/fig_split_big.pdf}
        \caption{Number of videos and splits}
        \label{fig:fun_split}
    \end{subfigure}
    \vspace{-25pt}
    \begin{subfigure}[b]{0.42\textwidth}
        \centering
        \includegraphics[width=1\textwidth]{figures/fig_vqa_num_big.pdf}
        \caption{Number of QA pairs}
        \label{fig:fun_qa}
    \end{subfigure}
    \begin{subfigure}[b]{0.5\textwidth}
        \centering
        \includegraphics[width=1\textwidth]{figures/fig_moment_distrib.png}
        \caption{Time Dist. of Localization}
        \label{fig:fun_moment}
    \end{subfigure}
    % \begin{subfigure}[b]{0.33\textwidth}
    %     \centering
    %     \includegraphics[width=1\textwidth]{figures/fig_avg_word_len.pdf}
    %     \caption{Avg. word count of answers}
    %     \label{fig:fun_avg}
    % \end{subfigure}
    \begin{subfigure}[b]{0.4\textwidth}
        \centering
        \includegraphics[width=1\textwidth]{figures/fig_val.png}
        \caption{Consensus Evaluation}
        \label{fig:fun_val}
    \end{subfigure}
    
    
  \vspace{-10pt}
  \caption{\textbf{Statistics of FunQA Dataset.} 
  % FunQA consists of three subsets, each corresponding to different video types, and is annotated with free-text QA pairs. 
  % The first row displays word clouds representing critical annotations for each subset. 
 Figure (a) and (b) showcase vital statistics, including the number of videos for different videos, splits, and QA pairs count for three subsets. Figure (c) highlights the high-frequency time span of the answer for localization questions in red. Figure (d) presents the percentage of consensus between annotators for the same QA pair. The consensus is categorized into three levels \textbf{High} for consistent understanding, \textbf{Medium} for partial agreement but with mutual acknowledgment, and \textbf{Low} for complete disagreement.
  }
  \vspace{-0.4cm}
  \label{fig:dataset statistics}
\end{figure}
% \begin{figure}[htbp]
%     \centering
%     \includegraphics[scale=0.5]{figures/dataset_main.pdf}
%     \caption{Pipleline of FunQA Dataset.}
%     \label{fig:dataset_main}
% \end{figure}

In this section, we provide a detailed explanation of the design principles that guided the creation of the FunQA dataset and its subsets. Additionally, we introduce our novel VideoQA tasks tailored for FunQA, and FunQA data statistics in Figure~\ref{fig:dataset statistics}. 
Toward the end we present our Construction Pipeline and Quality Control, highlighting our efforts in maintaining data quality and objectivity.


% \subsection{Video Selection}
% % \noindent\textbf{Design principle}\quad 
% In constructing the dataset, we adhered to three principles to address the challenges in video understanding capabilities: our dataset, FunQA, is \textbf{visual centered} and emphasizes \textbf{counter-intuitive reasoning}, \textbf{spatial-temporal reasoning}. Based on these principles, we collect 4365 videos from 3 different art genres and created three subsets: HumorQA, CreativeQA, and MagicQA.\\
% % \jingkang{Based on the principle, we collect xxx videos from 3 different art genres....... What is your subset, why these datasets, simple stats of these dataset.}
% \noindent\textbf{HumorQA}\quad
% HumorQA composed of 1,769 meticulously curated web videos, serves as a unique source of insight into human humor comprehension. Notably, it contains the shortest average video length of 7s among the three subsets. We believe that the human process of understanding humor is complex and deep, requiring a holistic understanding of the video and adding a degree of common sense to it. Psychological research has demonstrated that humor arises from the incongruity \cite{kant1987critique,sep-humor} between reality and expectations, flourishing with the skillful juxtaposition and transformation of events \cite{latta1999basic,boyd2004laughter,koestler2020act}. This makes humorous videos a valuable asset for the VideoQA dataset, anticipated to enhance a model's proficiency in integrating information and performing deep reasoning.\\
% \noindent\textbf{CreativeQA}\quad 
% CreativeQA is a collection of 927 videos averaging 48s in length from a TV show called Kasou Taishou \cite{ntv}. This program, showcasing original and novel skits performed by various amateur groups and judged by a panel, boasts a strong creative flair 
%  \cite{runco2012standard}. The essence of the show lies in using a mix of people and props to mimic reality, with audiences deriving pleasure from information integration and comparison. We anticipate that the imitation nature of the show will challenge the model's capacity for information extraction, while the longer video length and need for understanding creativity will put to test the model's comprehension of spatial-temporal information.\\
% \noindent\textbf{MagicQA}\quad
% MagicQA encapsulates 1672 magic performance videos sourced from across the web, spanning various genres like camera magic, close-up magic, and stage magic. The essence of magic revolves around the creation of seemingly impossible illusions \cite{nelms2012magic}, employing diverse effects such as disappearance, creation, and transformation. These illusions are infused with abundant spatial-temporal information. Through this dataset, we aim to empower the model to not only track the ensuing changes in objects but also unravel the underlying mechanics \cite{lamont2005magic} of these transformations.
% \noindent\textbf{FunQA}\quad
% In general, our meticulous filtering process ensures that the information within the aforementioned three types of art remains focused on visual content. These art forms naturally encompass complex spatial-temporal information, and their appreciation also requires deep human reasoning. Therefore, they align perfectly with our design guidelines.



\subsection{Task Definition}
% \jingkang{Each datasets has the following 4 types of tasks.....}


To comprehensively evaluate the model's ability to understand surprising videos, we designed the following 4 types of tasks for each subset: \\
\noindent\textbf{Counter-intuitive Timestamp Localization Task}\quad
The localization task is the base task to assess the model's comprehension abilities. It involves localizing counter-intuitive segments within the video, answers expressed in either seconds or frames. This task serves as the basis for the subsequent two main tasks in the three subsets, where the focus shifts to locating moments of humor, creativity, and magical effects, respectively. Successfully completing this task demands the model's understanding of the video's overall content, incorporating both temporal and spatial information.\\
\noindent\textbf{Detailed Description Task}\quad
The description task aims to evaluate the model's information extraction capabilities, serving as a fundamental aspect of video understanding. This task requires providing a free-text answer that describes the selected moment. Furthermore, this task allows for analysis of how the model extracts information and generates answers for subsequent tasks. By examining the model's performance in this task, we gain insights into its ability to extract relevant information and generate meaningful responses.\\
\noindent\textbf{Counter-intuitiveness Reasoning Task}\quad
The reasoning task is designed to test the model's ability to reason about the video, and in the three subsets, this question is Why Humorous, Why creative, and Why counter-intuitive and the answer is a free-text explanation. This task is very difficult and involves the model's deep reasoning ability; it requires the model to give a complete explanation using information from the entire video and its own common sense.\\
\noindent\textbf{Higher Level Tasks}\quad
In addition to the three main tasks, we design higher-level tasks to enhance the model's inference abilities on counter-intuitive videos. \textit{Title Task} in HumorQA and CreativeQA requires generating a concise title summarizing the video's content. \textit{Creative Scoring Task} in CreativeQA involves rating the creativity of videos between 1 and 20 (provided officially). \textit{Magic Method Task} in MagicQA requires the model to explain clearly the rationale behind the magic, and its purpose is to test the model's ability to reason more deeply. To ensure the accuracy of the answers, this task is only partially annotated and appears only in the test set, details of which can be found in Appendix \ref{appendix_a:dataset con}.

%In addition, some visual QA benchmarks \cite{Whoop} have begun to adopt free-text answers; this flexible answer format enables multimodal LLMs to go beyond multiple choice options (LLMs)~\cite{2023videochat,maaz2023videochatgpt}; for Video QA, however, free-text are less common. To bridge this gap, we also consider some \textbf{free-text answer} evaluations, demanding enhanced video comprehension capabilities from the model. It will also assess the model's ability to generate reasonably lengthy textual responses.
% \input{table/construction}

\subsection{Dataset Statistics}
% FunQA contains \textbf{4,365} counter-intuitive video clips and \textbf{311,950} question-answer pairs, the total length of these videos is \textbf{23.9} h and the average length of video clips is \textbf{19} seconds. The distribution of the video lengths can be seen in Figure \ref{fig:fun_len}. These videos were divided into train / val / test set and the exact quantities can be seen in Figure \ref{fig:fun_split}. FunQA consists of three fine-grained sub-datasets, each sub-dataset contains four well-designed tasks, the distribution of specific tasks and their corresponding QApairs can be seen in Figure \ref{fig:fun_qa}. Figure \ref{} shows the word cloud map of FunQA as a whole, while Figures \ref{fig:wc_h}, \ref{fig:wc_c} and \ref{fig:wc_m} show the word cloud maps of the three sub-data sets, respectively., the other detailed statistics of sub-datasets are given in Table \ref{T:main} and Appendix \ref{section: FunQA statics}.\\ 
% FunQA has a well-established annotation process and high annotation quality, our cross-validation statistics can be seen in Figure \ref{fig:fun_val}, and the average length of our answers is up to -words, the specific data can be seen in Figure \ref{fig:fun_avg}.Besides, the heat map of the distribution of the results of the Counter-intuitive Timestamp Localization Task over the whole video can be seen in Figure \ref{fig:fun_moment}, which shows that the labeled results are in line with common sense.
% % % static table----------------------------
\begin{table}[!t]
\centering
\caption{\textbf{Statistics of the FunQA dataset.}
\label{T:statics}}
\resizebox{0.8\linewidth}{!}{

\begin{tabular}{lccccccc}
\hline
\multirow{2}{*}{Datase} & \multicolumn{4}{c}{Videos} & \multirow{2}{*}{Avg len(s)} & \multirow{2}{*}{Total len(h)} & \multirow{2}{*}{\# QA Pairs} \\ \cline{2-5}
           & Train & Val & Test & Total &    &       &       \\ \hline
HumorQA    & 1,244  & 350 & 172  & 1,769  & 7  & 3.72  & 141,250  \\
CreativeQA & 657   & 185 & 85   & 927   & 48 & 12.53 & 78,740  \\
MagicQA    & 1,171  & 334 & 167  & 1,672  & 16 & 7.65  & 91,960  \\ \hline
FunQA      & 3,072  & 869 & 424  & 4,365  & 19 & 23.90 & 311,950 \\ \hline \\
\end{tabular}
}
\end{table}
%static table ----------------------------------
% \\
% The distribution of question word length is shown in \ref{fig:fun_avg}(a). Questions are on average 34.24 words, which is much longer than existing VideoQA datasets (e.g., 8.7 in Activity-QA \cite{yu2019activitynet}, 11.6 in NExT-QA \cite{xiao2021nextqanext} ). And the median of each subtask is slightly less than the mean. The distribution of Moment is shown in \ref{fig:fun_moment}(b). We can clearly find that the distribution of counter-intuitive emotions in three different videos is very different. Humor, Creative and Magic videos account for 75\%, 23\% and 61\% of the videos respectively. HumorQA, Creative, Magic QA word cloud is shown in \ref{fig:wc_h, fig:wc_h, fig:wc_m}(c), FunQA full word clouds are in \ref{Data Statistics}.
FunQA contains \textbf{4,365} counter-intuitive video clips and \textbf{311,950} QA pairs, the total length of these videos is \textbf{23.9h} and the average length of each clips is \textbf{19} seconds. FunQA consists three fine-grained subsets, each one containing well-designed tasks. The specific numbers of videos and splits can be seen in Fig. \ref{fig:dataset statistics} (a).
% Each subset's video lengths can be seen in Figure \ref{fig:dataset statistics} (c).
The specific number of QA pairs for each task can be seen in Fig.~\ref{fig:dataset statistics} (b). For our localization task, the heat map for the three different types of videos can be seen in Fig.~\ref{fig:dataset statistics} (c), which shows the high-frequency time span of the answer. 
% signifies that addressing the FunQA task requires a strong emphasis on video content, rather than relying solely on a few frames of images.
For the description and reasoning tasks, the average length of the words in their free-text answers reached \textbf{34.24}, which is much longer than existing VideoQA datasets (e.g., 2.6 in NExT-QA~\cite{xiao2021nextqanext}).
% The specific word count of each task is shown in Figure \ref{fig:dataset statistics} (d).
FunQA has a well-established annotation process and high annotation quality, the result of our annotation consensus evaluation are illustrated in Fig. \ref{fig:dataset statistics} (d).  Impressively, over 90\% of the annotations demonstrate a high level of consensus, while only 1\% exhibit low consensus. This clearly underscores the objectivity and reliability of the FunQA dataset.
% HumorQA, CreativeQA, and MagicQA word clouds are shown in Figure \ref{fig:dataset statistics} (a-c). More statistics and FunQA full word cloud are given in Appendix \ref{appendix_a: Dataset statics}.
\subsection{Dataset Construction Pipeline} 
\label{Section: dataset construction pipeline}
FunQA dataset construction pipeline was in three stages: Pre-processing, Manual Annotation, and Post-Processing. The whole process took about 900 hours with over 50 highly educated undergraduates as part-time annotators. See Appendix~\ref{appendix_a:dataset con} for more details on the dataset construction.\\
\noindent\textbf{Pre-Processing}\quad
Initially, we crawled videos from YouTube. Then we performed a two-stage manual cleaning and cutting process on the collection to ensure counter-intuitive features and video quality and to exclude non-ethical and sensitive content, resulting in video clips.\\
\noindent\textbf{Manual Annotation}\quad
We annotated the videos according to the characteristics of different task designs in Chinese. We screen and train the annotators to ensure the accuracy and high quality of the annotation, and finally produce the original annotated files. After the first round, we conducted a secondary round of 10\% of the tasks and performed Consensus Evaluation to ensure the objectivity.\\
\noindent\textbf{Post-Processing}\quad 
Based on our carefully designed tasks and high-quality annotations, we expanded our dataset using GPT-3.5. Firstly, we automatically translated the Chinese annotations into English. Subsequently, we generated more QA pairs that were faithful to the original ideas but presented differently. This not only made FunQA multilingual but also expanded its QA pair count to 312K. Additionally, we created diverse sub-dataset, FunQA-MC (multi-choice QA) and FunQA-DIA (dialogue QA). In addition, to focus on exploring the ability to handle counter-intuitive reasoning, we released FunQA-MC-R (a multi-choice version specifically containing counter-intuitive reasoning questions). More details are given in Appendix \ref{appendix_a: Augmentation} and Appendix \ref{appendix_b: FunQA Extended dataset}.
\begin{center}
\begin{minipage}{\textwidth}
\begin{minipage}[t]{0.48\textwidth}
\makeatletter\def\@captype{table}
\caption{\textbf{Consensus Evaluation Experiment.} This table shows the results from a random 10\% sample of QA pairs, cross-validated by annotators to assess agreement with existing annotations. `Low,' `Medium,' and `High' indicate the strength of the consensus.}
% \vspace{-5pt}
\resizebox{\linewidth}{!}{
\begin{tabular}{l|cccc}
\toprule
Consensus & HumorQA & CreativeQA & MagicQA & \# Total \\
% \toprule
\hline
Low & 1 & 1 & 2 & 4 \textbf{(1\%)} \\
Medium & 9 & 6 & 17 & 32 \textbf{(8\%}) \\
High & 199 & 111 & 194 & 504 \textbf{(91\%)}\\
\bottomrule
\end{tabular}}
\label{T:coneval}
\end{minipage}
\hfill
\begin{minipage}[t]{0.48\textwidth}
\makeatletter\def\@captype{table}
\caption{\textbf{Can FunQA be Solved Solely Based on Images?} The left two columns show the average number of questions answerable or not by humans using only 8 static, uniformly selected frames.}
\label{T:onlyframes}
% \vspace{-5pt}
\centering
%\setlength\tabcolsep{15pt}
\resizebox{\linewidth}{!}{
\begin{tabular}{l|ccc}
\toprule
Dataset  & \# Can & \# Cannot & Cannot Rate  \\
% \toprule
\hline
HumorQA & 7.8 & 32.2 & 80\% \\
CreativeQA & 6.1 & 33.9 & 85\% \\
MagicQA & 2.6 & 37.4 & 94\% \\
\midrule
\textbf{FunQA} & 16.5 & 103.5 & \textbf{86\%} \\
\bottomrule
\end{tabular}
}
\end{minipage}
\end{minipage}
\vspace{10pt}
\end{center}
\subsection{Quality Control}
\noindent\textbf{Minimal Errors and High Objectivity in FunQA}\quad
We assure that every annotation included in the final release of FunQA has been subjected to rigorous \textbf{multi-person}, \textbf{multi-round} review processes. Furthermore, We did manual consensus evaluation on released FunQA dataset, randomly sampling 10\% of the data from all three sub-datasets (HumorQA, CreativeQA, and MagicQA). As shown in Table \ref{T:coneval}, we get the 91\% high consensus.\\
\noindent\textbf{FunQA Emphasis on Temporal Dynamics}\quad
FunQA requires a strong emphases on temporal dynamics rather than solely on few frames of images. To prove that, we did the quantitative human experiments - We randomly selected 40 videos from each of the three sub-datasets, totaling 120 videos. For each video, we sampled 8 frames evenly. We enlisted 10 individuals who had not seen any FunQA videos before. We had them view the sequence of 8 consecutive frames and then watch the original video along with its annotations. They were asked to determine whether they can understand and answer the counter-intuitive understanding of the original video solely based on the images. Nearly \textbf{86\%} people thought that FunQA cannot be solved only by images, as shown in Table~\ref{T:onlyframes}.\\
%\begin{table}[!t]
\vspace{-0.05mm}
\caption{\textbf{Can FunQA be Solved Solely Based on Images?} The left two columns show the average number of questions answerable or not by humans using only 8 static, uniformly selected frames.}
\label{T:onlyframes}
\vspace{-10pt}
\centering
\setlength\tabcolsep{15pt}
\resizebox{\linewidth}{!}{
\begin{tabular}{l|ccc}
\toprule
Dataset  & \# Can & \# Cannot & Cannot Rate($\%$)  \\
% \toprule
\hline
HumorQA & 7.8 & 32.2 & 80 \\
CreativeQA & 6.1 & 33.9 & 85 \\
MagicQA & 2.6 & 37.4 & 94 \\
\midrule
\textbf{FunQA} & 16.5 & 103.5 & \textbf{86} \\
\bottomrule
\end{tabular}
}
\vspace{-5pt}
\end{table}


% \begin{table}[!t]
% \vspace{-0.05mm}
% \caption{\textbf{Can FunQA be Solved Solely Based on Images?} }
% \label{T:onlyframes}
% \vspace{-10pt}
% \centering
% \setlength\tabcolsep{15pt}
% \resizebox{0.8\linewidth}{!}{
% \begin{tabular}{l|cc}
% \toprule
% Dataset  & # Can & # Cannot & ($\%$)  \\
% % \toprule
% \hline
% HumorQA & 7.8 & 32.2 & 80 \\
% CreativeQA & 6.1 & 33.9 & 85 \\
% MagicQA & 2.6 & 37.4 & 94 \\
% \textbf{FunQA} & 16.5 & 103.5 & \textbf{86} \\
% \bottomrule
% \end{tabular}
% }
% \end{table} 
%\begin{table}[!t]
\begin{minipage}{0.5\linewidth}
\vspace{-0.05mm}
\caption{\textbf{Consensus Evaluation Experiment.} This table shows the results from a random 10\% sample of QA pairs, cross-validated by annotators to assess agreement with existing annotations. `Low,' `Medium,' and `High' indicate the strength of the consensus.}
\label{T:coneval}
\vspace{-10pt}
\centering
\setlength\tabcolsep{8pt}
\resizebox{\linewidth}{!}{
\begin{tabular}{l|cccc}
\toprule
Consensus & HumorQA & CreativeQA & MagicQA & \# Total \\
% \toprule
\hline
Low & 1 & 1 & 2 & 4 \textbf{(1\%)} \\
Medium & 9 & 6 & 17 & 32 \textbf{(8\%}) \\
High & 199 & 111 & 194 & 504 \textbf{(91\%)}\\
\bottomrule
\end{tabular}
}
\vspace{-10pt}
\end{minipage}%
\end{table}



% only frame
% consensus

% \begin{table}[!t]
\centering
\caption{\textbf{Comparison between FunQA and other existing benchmarks.}
Compared to other datasets, FunQA revolves around the captivating realm of interesting and counter-intuitive videos. The tasks within FunQA are specifically designed to challenge the vision capabilities of models, requiring strong skills in producing an in-depth description, interpretation, and spatial-temporal reasoning.
Here we clarify the abbreviation in the table. For annotation type:  \includegraphics[width=0.02\textwidth]{figures/human.png} denotes Manual Annotation and  \includegraphics[width=0.02\textwidth]{figures/robot.png} for Automatic Annotation; \textbf{Avg Len} denotes video average length; \textbf{\# Clips} means number of video clips; \textbf{VC} for visual-centric, \textbf{Des.} for Description, \textbf{Exp.} for Explanation, \textbf{STR} for Spatial-temporal Reasoning, \textbf{MC} means Multiple Choice QA, and \textbf{OE} shows Open Ended QA with \textbf{Average Word Count} per response. 
}
\vspace{-5pt}
\label{T:main}
\setlength\tabcolsep{10pt}
\resizebox{1.0\linewidth}{!}{
\begin{tabular}{llccccccccccl}
\toprule
\multirow{2}{*}{Dataset} & \multirow{2}{*}{Domain} & \multirow{2}{*}{
\includegraphics[width=0.02\textwidth]{figures/human.png} or \includegraphics[width=0.02\textwidth]{figures/robot.png}
} & \multicolumn{2}{c}{Video} & \multicolumn{7}{c}{Question Answer}  \\
\cmidrule(l){4-5} \cmidrule(l){6-12} 
 &  &  & Avg Len & \# Clips                      & \# QA  & VC & Des. & Exp. & STR & MC & OE   \\
 \toprule
TGIF-QA \cite{jang-IJCV-2019} 
& Social Media & \includegraphics[width=0.02\textwidth]{figures/robot.png} & 3s & 72K & 165K & \textcolor{darkgreen}{\ding{51}} & \textcolor{darkgreen}{\ding{51}} & \textcolor{darkred}{\ding{55}} & \textcolor{darkgreen}{\ding{51}} & \textcolor{darkred}{\ding{55}} & 2.1 \\

MSRVTT-QA \cite{xu2017video}
& Social Media & \includegraphics[width=0.02\textwidth]{figures/robot.png} & 15s & 10K & 244K & \textcolor{darkgreen}{\ding{51}} & \textcolor{darkred}{\ding{55}} & \textcolor{darkred}{\ding{55}} & \textcolor{darkgreen}{\ding{51}} & \textcolor{darkred}{\ding{55}} & 1.0 \\

ActivityNet-QA \cite{yu2019activitynetqa}
& Social Media & \includegraphics[width=0.02\textwidth]{figures/human.png} & 180s & 6K & 58K & \textcolor{darkgreen}{\ding{51}} & \textcolor{darkred}{\ding{55}} & \textcolor{darkred}{\ding{55}} & \textcolor{darkgreen}{\ding{51}} & \textcolor{darkred}{\ding{55}} & 1.9 \\
\toprule

NExT-QA \cite{xiao2021nextqanext}            
& Daily life & \includegraphics[width=0.02\textwidth]{figures/human.png} & 44s & 5K & 52K & \textcolor{darkgreen}{\ding{51}} & \textcolor{darkgreen}{\ding{51}} & \textcolor{darkgreen}{\ding{51}} & \textcolor{darkgreen}{\ding{51}} & \textcolor{darkgreen}{\ding{51}} & 2.6 \\

Social-IQ \cite{zadeh2019social}                        
& Daily life & \includegraphics[width=0.02\textwidth]{figures/human.png} & 99s & 1K & 8K & \textcolor{darkgreen}{\ding{51}} & \textcolor{darkred}{\ding{55}} & \textcolor{darkgreen}{\ding{51}} & \textcolor{darkred}{\ding{55}} & \textcolor{darkgreen}{\ding{51}} & N/A \\
\toprule

MovieQA \cite{tapaswi2016movieqa}                       & TV shows & \includegraphics[width=0.02\textwidth]{figures/robot.png} & 203s & 7K & 6K & \textcolor{darkred}{\ding{55}} & \textcolor{darkred}{\ding{55}} & \textcolor{darkgreen}{\ding{51}} & \textcolor{darkgreen}{\ding{51}} & \textcolor{darkgreen}{\ding{51}} & N/A \\

TVQA+ \cite{lei2019tvqa}  
& TV shows & \includegraphics[width=0.02\textwidth]{figures/human.png} & 8s & 4K & 30K & \textcolor{darkred}{\ding{55}} & \textcolor{darkred}{\ding{55}} & \textcolor{darkgreen}{\ding{51}} & \textcolor{darkgreen}{\ding{51}} & \textcolor{darkgreen}{\ding{51}} & N/A \\
\toprule


SUTD-TrafficQA \cite{xu2021sutd}
& Traffic & \includegraphics[width=0.02\textwidth]{figures/human.png} & 5s & 10K & 623K & \textcolor{darkgreen}{\ding{51}} & \textcolor{darkred}{\ding{55}} & \textcolor{darkred}{\ding{55}} & \textcolor{darkgreen}{\ding{51}} & \textcolor{darkgreen}{\ding{51}} & N/A \\
MarioQA \cite{mun2017marioqa} 
& Games  & \includegraphics[width=0.02\textwidth]{figures/human.png} & 5s & 188K & 188K & \textcolor{darkgreen}{\ding{51}} & \textcolor{darkred}{\ding{55}} & \textcolor{darkgreen}{\ding{51}} & \textcolor{darkgreen}{\ding{51}} & \textcolor{darkred}{\ding{55}} & 2.0 \\

CLEVRER \cite{yi2020clevrer}     
& Synthetic Videos & \includegraphics[width=0.02\textwidth]{figures/human.png} & 5s & 20K & 305K & \textcolor{darkgreen}{\ding{51}} & \textcolor{darkred}{\ding{55}} & \textcolor{darkgreen}{\ding{51}} & \textcolor{darkgreen}{\ding{51}} & \textcolor{darkgreen}{\ding{51}} & N/A \\


% \midrule

% HumorQA                                           
% & Daily life & M & 7 & 1.8 &  A & 141.3 & \textcolor{darkgreen}{\ding{51}} & \textcolor{darkgreen}{\ding{51}} & \textcolor{darkgreen}{\ding{51}} & \textcolor{darkgreen}{\ding{51}} & OE \& FT \\

% CreativeQA                                     
% & Performance & M & 48 & 0.9 &  A & 78.7 & \textcolor{darkgreen}{\ding{51}} & \textcolor{darkgreen}{\ding{51}} & \textcolor{darkgreen}{\ding{51}} & \textcolor{darkgreen}{\ding{51}} & OE \& FT \\

% MagicQA                                  
% & Magic & M & 10 & 1.6 &  A & 91.9 & \textcolor{darkgreen}{\ding{51}} & \textcolor{darkgreen}{\ding{51}} & \textcolor{darkgreen}{\ding{51}} & \textcolor{darkgreen}{\ding{51}} & OE \& FT \\
\midrule
\textbf{FunQA (Ours)} 
& \textbf{Surprising Videos} 
& \includegraphics[width=0.02\textwidth]{figures/human.png} & 19s & 4K & \textbf{312K} & \textcolor{darkgreen}{\ding{51}} & \textcolor{darkgreen}{\ding{51}} & \textcolor{darkgreen}{\ding{51}} & \textcolor{darkgreen}{\ding{51}} & \textcolor{darkgreen}{\ding{51}} & \textbf{34.2} \\ 
\bottomrule 

\end{tabular}}
\vspace{-0.5cm}
\end{table}


% \subsection{Dataset Comparision}

% As shown by the benchmarks in Table \ref{T:main}, there is growing interest in using question responses to measure video reasoning ability \cite{tapaswi2016movieqa, lei2018tvqa, lei2019tvqa, mun2017marioqa, jang-IJCV-2019, yi2020clevrer, xiao2021nextqanext, garcia2020knowit, grundemclaughlin2021agqa, wu2021star, Gao_2021_ICCV, castro2022fiber}. However, the data sources for these well-known benchmarks are often generated semi-automatically from well-known TV shows or films \cite{tapaswi2016movieqa, lei2018tvqa, lei2019tvqa}, so their tasks and answers often rely on dialogue and plot summaries rather than the visual content of the video \cite{tapaswi2016movieqa, lei2018tvqa, lei2019tvqa, kim2017deepstory, zadeh2019social}, leading to models that rely more on dialogue than on visual input, thus reducing the validity of the benchmarks for measuring visual reasoning \cite{tapaswi2016movieqa, lei2018tvqa, lei2019tvqa}.
% An increasing number of benchmarks are now often designed specifically to surprise by focusing on enhancing a particular domain of the model. \cite{yang2022avqa} focuses on improving the model's ability to process audio during problem solving. Moreover, \cite{xiao2021nextqanext, grundemclaughlin2021agqa, wu2021star}  currently focus on improving spatio-temporal reasoning.
% Other pure video benchmarks suffer from biases and simplicity associated with human-generated questions \cite{yu2019activitynetqa, tapaswi2016movieqa, jang-IJCV-2019, lei2018tvqa} or descriptions \cite{xu2017video, zeng2017leveraging}. The largest human annotated \cite{lei2018tvqa} and generated \cite{maharaj2017dataset} datasets contain 152.5K and 349K questions. Picture understanding has recently been increasingly studied in the context of strange, bizarre, and anomalous understanding, but there are no relevant datasets or benchmarks in the field of video understanding. FunQA surpasses the current VideoQA datasets in several aspects:\\
%  1) FunQA distinguishes itself through its diverse video types, each presenting unique characteristics within its respective subset. HumorQA involves incongruity and common sense elements, CreativeQA showcases intricate imitation behavior, and MagicQA entails the impossibility. FunQA's videos are longer and more informative, providing ample visual content for analysis.  \\
% 2) Secondly, FunQA places a strong emphasis on inference in its task design. This necessitates the model to possess deeper inference capabilities and leverage spatial and temporal information effectively. \\
% 3) Among the available VideoQA datasets, FunQA stands out for its inclusion of free-text type answers. Notably, the average answer length in Detailed Description Task and Counter-intuitiveness Reasoning Task is significantly longer compared to existing VideoQA datasets. This demands a higher output capability from the model and allows for evaluation using complex metrics. \\
% 4) FunQA benefits from high-quality human annotation, ensuring the reliability and accuracy of the dataset. They allow us to extend the FunQA dataset with GPT to suit different needs, e.g. FunQA Dialog is more in line with the LLM dialogue format, FunQA Multi-choice has proven evaluation metrics and testing is not limited to LLM, any model can be tested on it.


